\topic{完整案例推导样例}

下面六个案例不是为了增加题量，而是给全书提供一套可模仿的案例讲法。每个案例都按同一条线展开：先翻译题面，写暴力基线，再定位最贵的重复工作，随后设计状态、不变量和边界测试。以后读到题卡时，请用这一节的粒度要求自己，而不是只记“用什么算法”。

\subsection{LeetCode 76 最小覆盖子串：窗口从暴力检查中长出来}

题目给两个字符串 \code{s} 和 \code{t}，要求在 \code{s} 中找一个最短连续子串，使它包含 \code{t} 中每个字符及其次数。先把题意翻译成算法语言：答案对象是连续区间，合法性由字符频次决定，目标是在合法区间中找最短长度。暴力做法是枚举所有左端和右端，每得到一个子串就重新统计频次并判断是否覆盖 \code{t}。若字符串长度为 \code{n}、字符集大小为 \code{|\Sigma|}，这个做法至少是 \complexity{O(n^2|\Sigma|)}，如果每次重新切片或排序，还会更慢。

以 \code{s = "ADOBECODEBANC"}、\code{t = "ABC"} 为例，暴力从左端 0 开始会先找到 \code{"ADOBEC"}。它已经覆盖 A、B、C，继续向右扩张只会让当前左端对应的窗口更长，所以对“最短”没有帮助。于是可以尝试移动左端。移走 \code{A} 后窗口变成 \code{"DOBEC"}，A 不足，窗口不合法，必须继续右扩直到重新遇到 A。这个手算过程说明了窗口策略：右端负责让窗口获得覆盖能力，左端负责在已经覆盖时尽量缩短。

\begin{center}
\begin{tabular}{@{}p{0.18\linewidth}p{0.22\linewidth}p{0.44\linewidth}@{}}
阶段 & 当前窗口 & 结论\\
右端扩张到 C & \code{ADOBEC} & 第一次覆盖目标，记录长度 6。\\
左端收缩 & \code{DOBEC} & 移走 A 后不再覆盖，停止收缩。\\
继续扩张到 A & \code{DOBECODEBA} & 重新覆盖，左侧无关字符可以继续删除。\\
继续扩张到 C & \code{BANC} & 得到长度 4 的更短合法窗口。\\
\end{tabular}
\end{center}

现在再设计状态。若每次判断覆盖都全量比较计数表，仍然比暴力好，但还可以更清楚。 \code{need[ch]} 表示目标需要多少个字符 \code{ch}，\code{window[ch]} 表示当前窗口里有多少个。 \code{required} 是目标中不同字符种类数，\code{formed} 是当前已经满足需求的字符种类数。当某个字符从不足变成刚好满足，\code{formed} 加一；当左端移出字符导致它从满足变成不足，\code{formed} 减一。窗口合法条件就变成 \code{formed == required}。

不变量是：每次进入内层 \code{while} 时，当前窗口覆盖 \code{t}；内层循环每收缩一次都先记录合法答案，再移走左端字符并更新状态；一旦窗口不合法，停止收缩，交给右端继续扩张。每个字符最多进入窗口一次、离开窗口一次，所以时间是 \complexity{O(n + |\Sigma|)}。边界测试要包含 \code{t} 有重复字符、\code{s} 缺少目标字符、最短窗口在开头或结尾、大小写混合、\code{s} 比 \code{t} 短。这个案例的关键不是“用窗口”，而是证明连续区间合法性可以用增删字符维护。

\subsection{LeetCode 560 和为 K 的子数组：前缀哈希不是窗口替代品}

题目要求统计和为 \code{k} 的连续子数组数量，数组中可以有负数。先写暴力：枚举左端 \code{left}，再枚举右端 \code{right}，维护当前区间和，等于 \code{k} 就计数。这个版本覆盖所有连续子数组，复杂度 \complexity{O(n^2)}。它已经比每个区间重新求和的 \complexity{O(n^3)} 好，因为固定左端时复用了当前和；但每个右端仍然要回头看很多历史左端。

为什么不能直接用滑动窗口？因为负数会破坏单调性。右端加入一个负数，窗口和可能变小；左端移走一个负数，窗口和可能变大。于是“和太大就左移、和太小就右移”的判断不再可靠。真正可用的是前缀和代数：若 \code{prefix[j] - prefix[i] == k}，那么区间 \code{[i, j)} 的和就是 \code{k}。扫描到当前前缀 \code{prefix[j]} 时，只需要知道历史里有多少个 \code{prefix[j] - k}。

\begin{center}
\begin{tabular}{@{}p{0.12\linewidth}p{0.12\linewidth}p{0.14\linewidth}p{0.18\linewidth}p{0.30\linewidth}@{}}
下标 & 数值 & 当前前缀 & 需要历史 & 新增答案\\
初始 & -- & 0 & -- & 先记录空前缀一次。\\
0 & 1 & 1 & -2 & 没有。\\
1 & 2 & 3 & 0 & 找到 \code{[0,1]}。\\
2 & 1 & 4 & 1 & 找到 \code{[1,2]}。\\
3 & 2 & 6 & 3 & 找到 \code{[2,3]}。\\
\end{tabular}
\end{center}

哈希表的 key 是历史前缀和，value 是出现次数。为什么是次数？因为同一个前缀和可能出现多次，每一次都对应一个不同左端。为什么 \code{prefix_count[0] = 1}？它表示空前缀出现一次，让从下标 0 开始的合法区间能被统计到。为什么先查询再更新？当前前缀应该和历史前缀配对；如果先更新，在 \code{k = 0} 时会把空区间误算进去。

不变量是：处理当前元素之前，哈希表保存所有当前下标之前的前缀和出现次数；处理当前元素后，先用当前前缀查询答案，再把当前前缀加入历史。时间 \complexity{O(n)}，空间 \complexity{O(n)}。边界测试要包含负数、零、\code{k = 0}、从下标 0 开始的答案、多个相同前缀和、答案数量大于 1。这个案例的迁移点是 value 语义：问数量保存次数，问最长长度保存最早下标，问是否存在保存布尔或下标。

\subsection{LeetCode 875 爱吃香蕉的珂珂：答案二分先证明谓词}

题目给若干堆香蕉和小时数 \code{h}，要求找最小整数速度 \code{k}，使得能在 \code{h} 小时内吃完。先写暴力：速度从 1 开始试到最大堆大小，每个速度都计算总耗时，第一次可行就是答案。这个做法正确，因为答案一定在这个范围内；慢点是速度值域可能很大，每次只排除一个速度。

二分之前必须先定义谓词：\code{can_finish(speed)} 表示以这个速度能否在 \code{h} 小时内吃完。速度越大，每堆耗时 \code{ceil(pile / speed)} 越小或不变，所以如果某个速度可行，所有更大的速度也可行；如果某个速度不可行，所有更小的速度也不可行。答案空间呈现 false、false、true、true 的形状，我们要找第一个 true。

检查函数也要从题意推出来。每小时只能吃同一堆中的香蕉，一堆 \code{pile} 需要的小时数是向上取整的 \code{pile / speed}。整数写法是 \code{(pile + speed - 1) / speed}，但要注意溢出；稳妥做法是用 \code{std::int64_t} 计算。答案下界是 1，上界是最大堆大小，因为速度达到最大堆时，每堆最多一小时。

不变量可以使用左闭右开或左闭右闭，本书对“第一个可行值”常用闭区间收缩写法：\code{left} 和 \code{right} 都可能是答案；若 \code{mid} 可行，答案不在 \code{mid} 右侧，令 \code{right = mid}；若不可行，\code{mid} 及其左侧都不可行，令 \code{left = mid + 1}。循环结束时两者相等，就是最小可行速度。边界测试要包含只有一堆、\code{h} 等于堆数、速度需要等于最大堆、堆值接近 \code{int} 上限。这个案例的关键是：二分不是模板，真正的核心是谓词单调性和边界覆盖。

\subsection{LeetCode 84 柱状图中最大的矩形：弹栈时答案才确定}

题目给一排柱子，要求找最大矩形面积。先从暴力开始：枚举每根柱子作为矩形的最矮柱，向左找第一个比它矮的位置，向右找第一个比它矮的位置，宽度由这两个边界决定。这个做法正确，但每根柱子都可能向两边扫很远，最坏 \complexity{O(n^2)}。慢点在于“第一个更矮边界”被重复寻找。

单调栈的状态来自这个边界需求。维护一个下标栈，栈内柱高严格递增。扫描到新柱子 \code{i} 时，如果它比栈顶更矮，说明栈顶柱子的右侧第一个更矮位置已经找到，就是当前 \code{i}。弹出栈顶 \code{mid} 后，新的栈顶就是 \code{mid} 左侧第一个更矮位置；若栈空，说明左侧没有更矮位置。于是以 \code{heights[mid]} 为最矮柱的最大矩形可以结算。

\begin{center}
\begin{tabular}{@{}p{0.18\linewidth}p{0.24\linewidth}p{0.40\linewidth}@{}}
事件 & 栈含义 & 结算逻辑\\
新柱更高 & 递增关系保持 & 入栈，等待未来右边界。\\
新柱更矮 & 栈顶右边界确定 & 弹栈，当前下标是右侧更矮位置。\\
弹栈后 & 新栈顶是左边界 & 宽度为 \code{right-left-1}。\\
扫描结束 & 仍有柱子未结算 & 用哨兵高度 0 统一弹出。\\
\end{tabular}
\end{center}

为什么每根柱子只结算一次？因为它入栈时右侧更矮边界未知，弹出时右侧更矮边界第一次出现，之后它的最大矩形已经确定，未来不会再改变。每个下标最多入栈一次、出栈一次，所以时间 \complexity{O(n)}。相等高度可以用 \code{>=} 或 \code{>} 的不同策略处理，但必须保持宽度计算一致；常见写法是加一个尾部哨兵 0，让所有柱子在最后被弹出。边界测试要包含严格递增、严格递减、全相等、单柱、空数组、包含 0 的高度。

\subsection{LeetCode 994 腐烂的橘子：多源 BFS 是同时扩散}

题目给网格，腐烂橘子每分钟让四邻域的新鲜橘子腐烂，问全部腐烂最少分钟数。暴力模拟很自然：每分钟扫描整张网格，找旁边有腐烂橘子的新鲜橘子，把它们变烂。这个做法正确，但会反复扫描空格、已经腐烂的格子和本分钟不会变化的格子。

把题目翻译成图：每个新鲜或腐烂橘子格子是节点，四方向相邻是无权边，初始腐烂橘子都是源点。一个新鲜橘子的腐烂时间，就是它到最近源点的最短距离。因为所有边的代价都是一分钟，所以 BFS 第一次到达某个新鲜橘子时，就是它最早腐烂时间。多源 BFS 只是把所有源点一开始都放进队列，距离都设为 0。

队列层次就是时间。每轮外层循环开始时，队列里保存的是同一分钟会向外扩散的全部腐烂橘子。记录 \code{level_size}，只处理这一层；本层造成的新腐烂橘子入队，但要等下一分钟再扩散。每腐烂一个新鲜橘子就减少 \code{fresh_count}，最后若新鲜数为 0，分钟数就是答案；若队列空了还有新鲜橘子，说明它们不可达，返回 -1。

不变量是：队列中已经入队的格子都已经被标记为腐烂，避免同一格子被多个邻居重复入队；处理完第 \code{t} 层后，所有最短距离不超过 \code{t} 的可达新鲜橘子已经腐烂。边界测试要包含没有新鲜橘子、没有腐烂橘子、单个格子、被空格隔开的新鲜橘子、多个初始源点同时扩散。这个案例的迁移点是所有多源最短步数问题：先把所有源点以距离 0 入队，而不是为每个源点单独 BFS。

\subsection{LeetCode 72 编辑距离：DP 转移来自最后一步操作}

题目允许插入、删除、替换，把 \code{word1} 变成 \code{word2}，问最少操作数。暴力搜索会在每一步尝试三种操作，搜索树很快膨胀；更重要的是，不同操作序列会反复遇到同一个问题，例如“把 \code{word1} 的前 \code{i} 个字符变成 \code{word2} 的前 \code{j} 个字符”。这就是状态。

定义 \code{dp[i][j]}：把 \code{word1[0..i)} 变成 \code{word2[0..j)} 的最少操作数。初始化由空串决定：\code{dp[i][0] = i}，因为只能删除 \code{i} 次；\code{dp[0][j] = j}，因为只能插入 \code{j} 次。转移看最后一步。若末尾字符相同，最后一个字符不用动，继承 \code{dp[i-1][j-1]}。若不同，最后一步只有三种：删除 \code{word1} 末尾，来自 \code{dp[i-1][j] + 1}；插入 \code{word2} 末尾，来自 \code{dp[i][j-1] + 1}；替换末尾，来自 \code{dp[i-1][j-1] + 1}。

\begin{center}
\begin{tabular}{@{}p{0.16\linewidth}p{0.42\linewidth}p{0.24\linewidth}@{}}
操作 & 操作前已经解决的问题 & 候选值\\
删除 & 少一个 \code{word1} 字符时已能变成目标前缀 & \code{dp[i-1][j] + 1}\\
插入 & 当前 \code{word1} 已能变成少一个目标字符的前缀 & \code{dp[i][j-1] + 1}\\
替换 & 两边都少一个末尾字符时已经互相转换 & \code{dp[i-1][j-1] + 1}\\
\end{tabular}
\end{center}

遍历顺序按 \code{i}、\code{j} 递增即可，因为每个格子只依赖上方、左方和左上方。时间和空间都是 \complexity{O(mn)}。如果要压缩空间，必须保存左上角旧值，否则 \code{dp[j]} 被覆盖后会丢失替换来源。边界测试要包含空串、完全相同、完全不同、一个字符、前缀相同但后缀不同、后缀相同但前缀不同。这个案例的关键是：DP 公式不是背出来的，而是从最后一步操作逐项推出来的。

\begin{principlebox}{完整案例复盘要求}
读完一个案例后，请检查自己是否能回答六个问题：暴力枚举对象是什么，暴力为什么保证不漏答案，最贵的重复工作是哪一步，优化状态保存了什么，不变量如何证明不会漏答案，哪三个边界能打破错误写法。回答不出其中任何一个，就说明这题还只是“见过”，没有变成能力。
\end{principlebox}
